\subsection{MULTIPLIABLE FAMILIES IN C*}

In the multiplicative group \(\mathbf{C}^*\) of non-zero complex numbers a family \((z_{i})_{i\in I}\) cannot be multipliable unless \(\lim z_{i} = 1\) with respect to the filter of complements of finite subsets of \(I\) (Chapter III, § 5, no. 2, Proposition 1); furthermore, since every point of \(\mathbf{C}^*\) has a countable fundamental system of neighbourhoods, the set of indices \(i\) such that \(z_{i}\neq 1\) is countable if the family \((z_{i})\) is multipliable (Chapter III, § 5, no. 2, Corollary to Proposition 1).

PROPOSITION 2. A family \((z_{t})\) of complex numbers \(z_{t} = r_{t}(\cos \theta_{t} + i\sin \theta_{t})\) is multipliable if and only if the family \((r_{t})\) of absolute values of the \(z_{t}\) is multipliable in \(\mathbf{R}_{+}^{*}\) and the family \((\theta_{t})\) of amplitudes of the \(z_{t}\) is summable in the group of angles \(\mathfrak{A}\) .

In view of the structure of the group \(\mathbf{C}^*\) ( \(\S\) 1, no. 3, Proposition 1), the proposition is an immediate consequence of Proposition 4 of Chapter III, \(\S\) 5, no. 4.

If we map each angle \(\theta\) to that one of its measures (to any given base \(a\) ) which belongs to the interval \([- \frac{1}{2} a, \frac{1}{2} a]\) , we have a local isomorphism of \(\mathfrak{A}\) with \(\mathbf{R}\) ( \(\S 2\) , no. 2); since \(\lim \theta_{t} = 0\) with respect to the filter of complements of finite subsets of \(\mathbf{I}\) , we may replace the condition (in the statement of Proposition 2) that the family \(\theta_t\) should be summable in \(\mathfrak{A}\) by the condition that the family \((t_t)\) of measures of the angles \(\theta_t\) which belong to \([- \frac{1}{2} a, \frac{1}{2} a]\) should be summable in \(\mathbf{R}\) .

The following theorem gives another criterion for a family of complex numbers, put in the form \((\mathbf{1} + u_{\iota})\) , to be multipliable in \(\mathbf{C}^*\) . (It generalizes Theorem 4 of Chapter IV, § 7, no. 4; see also Chapter IX, Appendix, no. 2, Proposition 1):

THEOREM 1. The family \((\mathbf{1} + u_{\iota})_{\iota \in \mathbf{I}}\) is multipliable in \(\mathbf{C}^*\) if and only if the family \((|u_{\iota}|)\) is summable in \(\mathbf{R}\) .

For each finite subset J of I, put

\[
p _ {J} = \prod_ {\iota \in J} (\mathbf{1} + a _ {\iota}), \quad s _ {J} = \sum_ {\iota \in J} a _ {\iota}, \quad \sigma_ {J} = \sum_ {\iota \in J} | a _ {\iota} |.
\]

LEMMA 1. For each finite subset J of I, let \(\varphi(J) = \sup_{L \subset J} (p_L - 1)\) . Then for each subset L of J we have

\[
\left| p _ {\mathbf {L}} - \mathbf {1} - s _ {\mathbf {L}} \right| \leqslant \varphi (\mathrm{J}) \sigma_ {\mathbf {L}}.\tag{2}
\]

This is clear if L is empty. We proceed by induction on card (L). Let \(L = K \cup \{\lambda\}\) , where \(\lambda \notin K\) ; then \(p_{\mathbf{L}} = p_{\mathbf{K}}(\mathbf{1} + a_{\lambda})\) and \(s_{\mathbf{L}} = s_{\mathbf{K}} + a_{\lambda}\), so that \(p_{\mathbf{L}} - \mathbf{1} - s_{\mathbf{L}} = (p_{\mathbf{K}} - \mathbf{1} - s_{\mathbf{K}}) + (p_{\mathbf{K}} - \mathbf{1})a_{\lambda}\) ; hence, by the inductive hypothesis and the definition of \(\varphi(\mathbf{J})\) , we have

\[
\left| p _ {\mathrm{L}} - \mathrm{1} - s _ {\mathrm{L}} \right| \leqslant \varphi (\mathrm{J}) \sigma_ {\mathrm{K}} + \varphi (\mathrm{J}) \left| a _ {\lambda} \right| = \varphi (\mathrm{J}) \sigma_ {\mathrm{L}},
\]

which proves the lemma.

LEMMA 2. If J is a finite subset of I such that \(\varphi(J) < 1/4\) , then

\[
\left| \sigma_ {J} \right| \leqslant 4 \varphi (J) / (1 - 4 \varphi (J)).
\]

For since \(\sigma_{\mathbf{L}} \leqslant \sigma_{\mathbf{J}}\) for all subsets \(\mathbf{L}\) of \(\mathbf{J}\) , it follows from (2) that \(|s_{\mathbf{L}}| \leqslant \varphi(\mathbf{J})\sigma_{\mathbf{J}} + |p_{\mathbf{L}} - \mathbf{1}| \leqslant (\mathbf{1} + \sigma_{\mathbf{J}})\varphi(\mathbf{J})\) ; but by virtue of Chapter VII, §3, no. 1, Proposition 2, we have \(|\sigma_{\mathbf{J}}| \leqslant 4\sup_{\mathbf{L} \subset \mathbf{J}} |s_{\mathbf{L}}|\) , hence \(\sigma_{\mathbf{J}} \leqslant 4\varphi(\mathbf{J})(\mathbf{1} + \sigma_{\mathbf{J}})\) , and the result follows.

Now let us show that the condition stated in Theorem 1 is sufficient. The hypothesis that the family \((|u_{\iota}|)\) is summable in \(\mathbf{R}\) implies that the family \((\mathrm{1} + |u_{\iota}|)\) is multipliable in \(\mathbf{R}_{+}^{*}\) (Chapter IV, § 7, no. 4, Theorem 4); hence, for each \(\varepsilon > 0\) , there is a finite subset \(J_{0}\) of I such that, for each finite subset L of I which does not meet \(J_{0}\) , we have \(\prod_{\iota \in L} (\mathrm{1} + |u_{\iota}|) - \mathrm{1} \leqslant \varepsilon\) . But we can write \(\prod_{\iota \in L} (\mathrm{1} + u_{\iota}) - \mathrm{1}\) in the form \(\sum_{M} \left( \prod_{\iota \in M} u_{\iota} \right)\) where M runs through all non-empty subsets of L; and since \(\left| \prod_{\iota \in M} u_{\iota} \right| = \prod_{\iota \in M} |u_{\iota}|\) , we have

\[
\left| \prod_ {\iota \in L} (1 + u _ {\iota}) - 1 \right| \leqslant \sum_ {M} \left(\prod_ {\iota \in M} | u _ {\iota} |\right) = \prod_ {\iota \in L} (1 + | u _ {\iota} |) - 1 \leqslant \varepsilon .
\]

This proves our assertion, by virtue of Cauchy's criterion, since \(\mathbf{C}^*\) is a complete group.

We still have to show that the condition of Theorem 1 is necessary. If \((\mathbf{1} + u_{\iota})_{\iota\in \mathbf{I}}\) is a multipliable family in \(\mathbf{C}^*\) , there exists a finite subset \(J\) of \(I\) such that, for every finite subset \(H\) of \(I\) which does not meet \(J\) , we have \(\left|\prod_{\iota\in H}(1 + u_{\iota}) - 1\right| \leqslant 1/8\) . By Lemma 2, it follows that \(\sum_{\iota\in H}|u_{\iota}| \leqslant 1\) for every finite subset \(H\) of \(I\) which does not meet \(J\) , and hence the family \((|u_{\iota}|)\) is summable in \(\mathbf{R}\) (Chapter IV, § 7, no. 1, Theorem 1).

The proof above applies also, mutatis mutandis, to (ordered) infinite products in certain non-commutative division rings and algebras (see Exercise 6, and Chapter IX, Appendix).

